\section{Introduction}
\label{sec:intro}

Language-model agents are moving from research prototypes into deployed products, and one design question now precedes almost every deployment: should a task be handled by a single-agent system (SAS) or by a multi-agent system (MAS) whose agents plan, check and specialize? Multi-agent frameworks have multiplied, and industrial reports describe large gains from orchestrating several agents, at many times the tokens of a single one \citep{anthropic2025multiagent}. Studies that hold compute fixed \citep{tran2026single} or account for cost after the fact \citep{jwalapuram2026illusion} often find the SAS as strong or stronger. The choice is consequential, since a MAS multiplies model calls, latency and points of failure.

\citet{kim2026capable} gave the most systematic answer to date, and with it the beginnings of a science of agent scaling. Across 260 configurations spanning six benchmarks, five architectures and nine models, the capability of the underlying SAS was the strongest predictor of what coordination adds: a MAS pays while the SAS is weak and costs once it is strong, with the fitted crossing near 45\% single-agent accuracy. In a validation on 16 model-by-benchmark cells the rule predicted in 15 the direction of the gain of the best of its four MAS architectures over the SAS. The authors offer it as a practical selection rule derived on the tasks and architectures they tested rather than a universal scaling law; we call it the capability-saturation principle for short, since what we test is whether the rule, as practitioners would apply it, carries over to a new setting.

Every benchmark behind this principle is autonomous: the agent receives a task and acts alone. Many agents now reaching users operate in a different regime. The agent converses with a person who supplies information late, changes goals and asks for what the policy forbids, and in the dual-control case of \citet{barres2025tau2} the person holds tools of their own and acts on the shared environment too. This regime bears on a MAS in three ways. The conversation has one user-facing locus, so most generic MAS architectures have no deployable analogue. Input tokens dominate the cost, so what is held constant when comparing systems decides the outcome. And the user is simulated by a model whose choice alone can move scores as much as the per-cell effects reported for the principle \citep{seshadri2026lost,zhou2026sim2real}. Whether principles derived without a user hold once a user takes part is of both scientific and practical weight.

We present such a test on $\tau^2$-bench \citep{barres2025tau2}, the standard dual-control benchmark, where the reward is the state of a database or device when the conversation ends; its airline and retail domains give tools to the agent alone, and its telecom domain gives the user tools as well. We compare the benchmark's stock SAS with a centralized MAS built from the same model: an orchestrator that plans and talks to the customer, a database agent that executes its reads and writes, and a policy advisor (Figure~\ref{fig:design}). Six models from three families on three domains give 18 cells, of which 16 were run on the registered configuration, for about 7,950 simulated conversations in all. The study was pre-registered, and every change to the plan is dated in a public change log.\footnote{Code, prompts, the pre-registered plan with its change log, all trajectories with per-call token records, budget ledgers and calibration files are available at \todo{public repository URL}.} Within a cell both systems run the same 50 tasks with three seeds, the MAS under a per-episode completion-token ceiling calibrated from the SAS's own spend. We read the result in one of three ways: the principle \emph{transfers} if the gain declines with capability and crosses zero near 45\%, \emph{shifts} if it crosses elsewhere, and \emph{disappears} if no dependence is detected. Mechanical validity checks separate a MAS that does not help from one that is broken, and the primary test uses a capability index measured on held-out domains, which removes the negative slope that a gain regressed on its own baseline shows by construction.

We find no support for the principle's distinctive prediction in this setting, for this architecture. Under a fixed simulated user no cell shows a MAS benefit: the fitted gain is negative across the observed capability range, with a weighted mean of $-28$ points, and the MAS is worse than the SAS in 12 of the 13 cells that meet the mechanical validity bars and tied in the thirteenth, the cell nearest 45\%, where only a change of simulator produces a 10-point gain (\S\ref{sec:robust}). On the two cells where the principle predicts a benefit, the MAS loses by 13 and 22 points, and with no budget ceiling at all it at best draws level at 5.5 times the SAS's median token spend (\S\ref{sec:transfer}). Whether the gain declines with capability the design could not resolve: the slope test detects no dependence, and its interval cannot exclude a steep decline, because model-by-domain effects dominate the variation between cells (\S\ref{sec:h1}). The mechanism that motivates a MAS in this setting, intercepting wrongful database writes, is not detectable: the MAS does not measurably reduce them and converts other failures into budget exhaustion and inaction (\S\ref{sec:h4}). One qualification runs through the paper. The MAS as built did not meet one of our pre-registered validity bars: its orchestrator planned without the tool schemas that only the database agent holds, and about three quarters of its write orders were not executed exactly as written. The pre-registered response, a fix and a re-run, was built and smoke-tested but not run on the grid; we report the architecture as built and label the deviation (\S\ref{sec:compliance}).

\paragraph{Contributions.} (1) The first test of the capability-saturation principle in the regime where agents serve users in conversation, on the dual-control benchmark. (2) A compute-matching protocol for multi-turn agent systems: one per-episode completion-token ledger for all agents of a system, calibrated per cell. (3) An analysis, supported by a lemma, that separates capability dependence from the coupling between a gain and its own baseline. (4) Evidence on why the MAS fails: a taxonomy separating write-type from word-type violations, per-cell validity checks, and a no-ceiling arm that rules out the budget as the sole cause. (5) Release of all trajectories, ledgers, prompts and code (footnote~1).
